% ============================================================================
% MUSTERLÖSUNG DS 6 - Gruppe A + B: Personas y características
% ============================================================================

\documentclass[11pt,a4paper]{article}

\usepackage[utf8]{inputenc}
\usepackage[T1]{fontenc}
\usepackage[ngerman]{babel}
\usepackage[left=2cm,right=2cm,top=1.8cm,bottom=1.5cm]{geometry}
\usepackage{helvet}
\usepackage{xcolor}
\usepackage{tabularx}
\usepackage{array}
\usepackage{enumitem}
\usepackage{tcolorbox}
\usepackage{fancyhdr}
\usepackage{lastpage}

\renewcommand{\familydefault}{\sfdefault}

\definecolor{spanisch}{HTML}{c41e3a}
\definecolor{grau}{HTML}{666666}
\definecolor{hellgrau}{HTML}{f5f5f5}
\definecolor{gruppeB}{HTML}{1565c0}
\definecolor{loesung}{HTML}{c62828}

\pagestyle{fancy}
\fancyhf{}
\fancyhead[L]{\small\textcolor{grau}{Spanisch Spätbeginner -- MUSTERLÖSUNG}}
\fancyhead[R]{\small\textcolor{grau}{ML-06}}
\fancyfoot[C]{\small\thepage/\pageref{LastPage}}
\renewcommand{\headrulewidth}{0.5pt}

\newcommand{\lsg}[1]{\textcolor{loesung}{\textbf{#1}}}
\newcommand{\es}[1]{\textit{#1}}

\newtcolorbox{gruppenbox}[1][]{
    colback=white,
    colframe=spanisch,
    fonttitle=\bfseries\color{white},
    colbacktitle=spanisch,
    title=#1,
    boxrule=1pt,
    arc=0pt,
    left=5pt, right=5pt, top=5pt, bottom=5pt
}

\newtcolorbox{gruppeBbox}[1][]{
    colback=white,
    colframe=gruppeB,
    fonttitle=\bfseries\color{white},
    colbacktitle=gruppeB,
    title=#1,
    boxrule=1pt,
    arc=0pt,
    left=5pt, right=5pt, top=5pt, bottom=5pt
}

\begin{document}

% ============================================================================
% KOPFBEREICH
% ============================================================================
\begin{center}
{\Large\textbf{\textcolor{loesung}{MUSTERLÖSUNG}}}\\[0.3em]
{\large DS 6 -- Brückenthema: Personas y características}
\end{center}

\vspace{0.5em}
\hrule
\vspace{1em}

% ============================================================================
% GRUPPE A
% ============================================================================
\begin{gruppenbox}[Gruppe A: Personenbeschreibung (33 Punkte)]

\textbf{Kasten 1: ser vs. estar + Adjektiv}

\begin{tabularx}{\textwidth}{@{}|p{2.5cm}|X|X|@{}}
\hline
\textbf{Verb} & \textbf{Verwendung} & \textbf{Beispiel} \\
\hline
\textbf{ser} & \lsg{Eigenschaften} (dauerhaft) & Mi hermana \lsg{es} simpática. \\
\hline
\textbf{estar} & \lsg{Zustände / Gefühle} (momentan) & Hoy \lsg{estoy} cansado. \\
\hline
\end{tabularx}

\vspace{0.5em}
\textbf{Kasten 2: Konnektoren} -- Beispiele für Lücken: \lsg{inteligente, paciente, perezoso, divertida} (variabel)

\vspace{0.8em}

\textbf{Aufgabe 1: Adjektive -- Endungen und Bedeutungen (8 P.)}

\begin{tabularx}{\textwidth}{@{}|X|X|X|X|@{}}
\hline
simpátic\lsg{o/-a} & \lsg{nett, sympathisch} & paciente & \lsg{geduldig} \\
\hline
inteligente & \lsg{intelligent, klug} & amable & \lsg{freundlich, liebenswürdig} \\
\hline
trabajador\lsg{/-a} & \lsg{fleißig, arbeitsam} & divertid\lsg{o/-a} & \lsg{lustig, unterhaltsam} \\
\hline
responsable & \lsg{verantwortungsvoll} & creativ\lsg{o/-a} & \lsg{kreativ} \\
\hline
\end{tabularx}

\vspace{0.5em}
\textbf{Aufgabe 2: ser oder estar (6 P.)}
\begin{enumerate}[label=\alph*), leftmargin=2em, nosep]
    \item Mi madre \lsg{es} muy paciente.
    \item Hoy mis hermanos \lsg{están} cansados.
    \item ¿Cómo \lsg{es} tu mejor amigo?
    \item Nosotros \lsg{somos} trabajadores.
    \item María \lsg{está} enferma hoy.
    \item Yo \lsg{soy} responsable y organizada.
\end{enumerate}

\vspace{0.5em}
\textbf{Aufgabe 3: Personenbeschreibung (8 P.) -- Beispiellösung:}

\lsg{Mi mejor amigo se llama Pablo. Es muy simpático y también es inteligente. Además, es muy divertido y siempre me hace reír. Pero a veces es un poco impaciente cuando tiene que esperar.}

\vspace{0.5em}
\textbf{Aufgabe 4: Partnerdialog (6 P.)} -- individuelle Antworten, z.B.:
\begin{enumerate}[leftmargin=2em, nosep]
    \item \lsg{[Name] es simpático/-a y trabajador/-a.}
    \item \lsg{Su mejor amigo/-a es muy divertido/-a.}
    \item \lsg{Su familia es grande y muy unida.}
\end{enumerate}

\vspace{0.5em}
\textbf{Aufgabe 5: Zusammenfassung (5 P.) -- Beispiel:}

\lsg{María es muy simpática y trabajadora. También es responsable. Su mejor amigo es Luis y es muy divertido.}

\end{gruppenbox}

\vspace{1em}

% ============================================================================
% GRUPPE B
% ============================================================================
\begin{gruppeBbox}[Gruppe B: Eigenschaften für Berufe (33 Punkte)]

\textbf{Kasten 1: hay que + Infinitiv}

Lücke: Hay que \lsg{trabajar} mucho.

\vspace{0.8em}

\textbf{Aufgabe 1: Berufe und Eigenschaften (8 P.)}

\begin{tabularx}{\textwidth}{@{}|X|X|X|@{}}
\hline
\textbf{Beruf} & \textbf{Deutsch} & \textbf{Eigenschaft} \\
\hline
el/la médico/-a & \lsg{der Arzt / die Ärztin} & \lsg{paciente, responsable} \\
\hline
el/la profesor/-a & \lsg{der Lehrer / die Lehrerin} & \lsg{paciente, organizado/-a} \\
\hline
el/la enfermero/-a & \lsg{der Krankenpfleger / die Krankenschwester} & \lsg{amable, paciente} \\
\hline
el/la jefe/-a & \lsg{der Chef / die Chefin} & \lsg{responsable, organizado/-a} \\
\hline
\end{tabularx}

\vspace{0.5em}
\textbf{Aufgabe 2: Sätze mit hay que ser ... porque ... (8 P.)}

\begin{enumerate}[label=\alph*), leftmargin=2em, nosep]
    \item Para ser profesor hay que ser \lsg{paciente} porque \lsg{trabajas con muchos niños / tienes que explicar muchas veces}.
    \item Para ser enfermero hay que ser \lsg{amable} porque \lsg{cuidas a personas enfermas / necesitan apoyo}.
    \item Para ser jefe hay que ser \lsg{responsable} porque \lsg{tomas decisiones importantes / diriges a un equipo}.
    \item Para ser \lsg{bombero} hay que ser \lsg{valiente} porque \lsg{el trabajo es peligroso}. (Beispiel)
\end{enumerate}

\vspace{0.5em}
\textbf{Aufgabe 3: ¿Cómo es un buen jefe? (6 P.) -- Beispiellösung:}

\begin{enumerate}[leftmargin=2em, nosep]
    \item \lsg{Un buen jefe es responsable y organizado.}
    \item \lsg{Además, tiene que ser paciente con sus empleados.}
    \item \lsg{Pero no es egoísta, escucha las opiniones de su equipo.}
\end{enumerate}

\vspace{0.5em}
\textbf{Aufgabe 4: Interview-Fragen (6 P.)}

\begin{enumerate}[label=\alph*), leftmargin=2em, nosep]
    \item Wie bist du? $\rightarrow$ \lsg{¿Cómo eres?}
    \item Wie ist dein bester Freund? $\rightarrow$ \lsg{¿Cómo es tu mejor amigo/-a?}
    \item Wie ist deine Familie? $\rightarrow$ \lsg{¿Cómo es tu familia?}
    \item Warum? $\rightarrow$ \lsg{¿Por qué?}
\end{enumerate}

\vspace{0.5em}
\textbf{Aufgabe 5: Notizvorlage und Zusammenfassung (5 P.)}

Individuelle Antworten basierend auf dem Interview. Beispiel für Zusammenfassung:

\lsg{Laura es simpática y muy divertida. También es inteligente. Su mejor amigo es muy trabajador.}

\end{gruppeBbox}

\vspace{1em}
\hrule
\vspace{0.3em}

\textbf{Bewertungshinweise:}
\begin{itemize}[leftmargin=1.5em, nosep]
    \item Gruppe A: 33 Punkte gesamt
    \item Gruppe B: 33 Punkte gesamt
    \item Bei Adjektiven: Endung (-o/-a) muss zum Subjekt passen
    \item Bei ser/estar: Begründung ist entscheidend (Eigenschaft vs. Zustand)
    \item Bei freien Texten: Inhalt wichtiger als perfekte Grammatik, aber Konnektoren müssen verwendet werden
    \item Halbe Punkte möglich bei teilweise richtigen Antworten
    \item Interview-Aufgaben: Flüssigkeit und Verständlichkeit bewerten
\end{itemize}

\vspace{0.5em}
\textbf{Typische Fehler -- Hinweise für Korrektur:}
\begin{itemize}[leftmargin=1.5em, nosep]
    \item \textcolor{loesung}{*Estoy inteligente} $\rightarrow$ \lsg{Soy inteligente} (Eigenschaft = ser)
    \item \textcolor{loesung}{*Mi hermano es simpático-a} $\rightarrow$ \lsg{Mi hermano es simpático} (maskulin!)
    \item \textcolor{loesung}{*Hay que es paciente} $\rightarrow$ \lsg{Hay que ser paciente} (Infinitiv!)
\end{itemize}

\end{document}
